% Source: https://artofproblemsolving.com/community/c3h2096167_number_theory_fun
% Stable identity: aops-topic-2096167-q03
Write $n=2^e\,p_1^{k_1}\cdots p_r^{k_r}$ with odd primes $p_i$. Then
\[\sigma(n)=\left(1+2+\cdots+2^e\right)\prod_{i}\left(1+p_i+\cdots+p_i^{k_i}\right).\]
The first factor is odd. The factor $1+p_i+\cdots+p_i^{k_i}$ is a sum of $k_i+1$ odd numbers, so it is odd exactly when $k_i$ is even. Hence $\sigma(n)$ is odd exactly when every odd prime appears to an even power, that is, when $n=m^2$ or $n=2m^2$ for some integer $m$.

Between $100$ and $200$: the squares are $100,121,144,169,196$ ($5$ numbers), and the doubled squares are $2\cdot8^2=128$, $2\cdot9^2=162$, $2\cdot10^2=200$ ($3$ numbers). So the answer is $\boxed{8}$.
